% ======================================================================
%  EXPERIMENTAL SETUP
% ======================================================================
\section{Experimental Setup}\label{sec:setup}

\subsection{Data Collection}

We collect live microstructure data from centralized cryptocurrency exchanges using a custom collector that polls 9~signal families (orderbook, trades, ticker, funding, open interest, withdrawal status, exchange status, spread matrices, OHLCV) via REST API at approximately 60-second target cadence across 23~assets and 6~exchanges (Binance, Bybit, OKX, Coinbase, Kraken, MEXC). Under full signal load, the observed cadence is approximately 110~seconds per snapshot. Data collection protocol and signal schema are documented in the repository accompanying \datasetref.

\subsection{Train/Test Split}

All collection windows before July~25, 2026 form the training set; July~25--28 forms the held-out test set. The split is strictly chronological with no shuffling. The training set contains approximately 2.9~million (coin, exchange-pair, snapshot) feature rows spanning multiple collection windows from June--July 2026. Cross-sectional rows within a snapshot are structurally dependent (they share the same market state), so this row count should not be interpreted as independent samples.

\subsection{Live Paper-Trading Protocol}

For live validation, we run independent paper-trading campaigns that process incoming collector data in real time with no access to future snapshots. Capacity is $\text{max\_open} = 50$. The paper confidence gate is $\tau{=}0.9$ (Section~\ref{sec:ablation}).

\subsubsection{Campaign A: Jul~31 - TO REMOVE}

\begin{itemize}
    \item \textbf{Model}: 68-feature prediction model (\texttt{statarb\_lgbm.txt}).
    \item \textbf{Horizon}: $H = 1$ (predict next snapshot).
    \item \textbf{Z-score window}: 300~snapshots.
    \item \textbf{Duration}: approximately 8~hours of useful data (collector ran for 264~snapshots).
    \item \textbf{Warmup}: predictions become active after $\text{MIN\_PERIODS} = 90$ snapshots (approximately 2.5--3~hours).
    \item \textbf{Predictions scored}: 58,995 total; 2,026 rows with $|\hat{z}|\geq 0.9$.
\end{itemize}

We validate reuslts on a 72\-hour live trading campaign, using the same model and model parameters as in training and testing. 

\subsubsection{Campaign B: Aug~4--7 (\textasciitilde72h)}

\begin{itemize}
    \item \textbf{Model}: same 68-feature, $H{=}1$, $w{=}300$ policy as Campaign~A.
    \item \textbf{Duration}: $\approx$72.9~hours of active trading (50,690 closed trades; 1.01M scored predictions).
    \item \textbf{Role}: multi-day live stress test; primary $\tau{=}0.9$ metrics are the post-hoc paper-protocol subset of this book (Section~\ref{sec:ablation}).
\end{itemize}

\subsection{Baseline Definitions - NEED TO UPDATE }

\subsubsection{Naive persistence}

On every row where the model generates a prediction, we also record the naive forecast $\hat{z}_{t+H} \leftarrow z_t$ (the current z-score predicts the future z-score unchanged). R$^2$ and directional accuracy are computed for both model and naive on the \emph{identical} rows, ensuring that any performance difference is attributable to the model rather than to row selection.

\subsubsection{Mechanical z-score baselines}

We replay two mechanical trading rules on the Jul~31 signal panel under the same capacity constraints ($\text{max\_open} = 50$), settlement proxy, and paper gate $\tau{=}0.9$ as the model:

\begin{itemize}
    \item \textbf{Persistence}: enter when $|z_t| \geq 0.9$, direction $= \text{sign}(z_t)$.
    \item \textbf{Mean-reversion}: enter when $|z_t| \geq 0.9$, direction $= -\text{sign}(z_t)$.
\end{itemize}

Scarce slots are filled by ranking $|\hat{z}|$ (model) or $|z_t|$ (mechanical). These represent the two canonical directions a rule-based z-score system can take: betting the spread continues (persistence) or reverts (classical pairs trading). Both settle at $t+1$ with $\text{pnl\_proxy} = \text{direction} \times z_{t+1}$, identical to the model's settlement.

\subsection{Reproducibility}

All training, evaluation, and baseline scripts are available in the code repository (\coderef), and the dataset is released at \datasetref. 